% Magda Gregorova, 2025-01-10
% abstract

\section*{Abstract (en)}
Current generative models can produce visually high-quality raster images from input
text or image, but the field falls short in vector graphics format due to its complex
nature. An SVG file is a document that contains paths, layers, and hierarchy and it
means its purpose is not just to represent something semantically, it should be structurally clean and neat. 
Those requirements create far more problems to overcome to
build a model for SVG generation. This thesis analyze SVG generation methods around fourteen open-source model 
to understand how internal design choices affect the structure of output. The analysis is shaped around three dimensions, 
how models represent geometry of vector internally, what supervision signals lead the generation process and which metrics 
are useful to evaluate the SVG. Representation and supervision interact rather than operate independently and they are two 
main parts of internal design of models. The analysis showed that same loss function can produce different structural outcomes 
depending on different representation choices and the same representation behaves differently under different supervision. 
The most consistent pattern across the models is a trade-off between visual realism and editability. Models that generate the 
most visually detailed outputs also produce the most fragmented SVG files, although methods that constrain their representation 
or learn from data produce cleaner results but within narrower domains. Current evaluation metrics (FID, CLIP, LPIPS) all operate 
on rendered pixels and cannot detect structural differences. To address this, the thesis proposes a six-dimension evaluation 
framework that includes path economy, geometric quality, layer organization, and editability alongside visual metrics. 
The thesis does not propose a new model. It provides a framework for understanding why current models produce the structures 
they do and what needs to change for generated SVGs to work in real design tools.


\section*{Abstract (de)}
Aktuelle generative Modelle können visuell hochwertige Rasterbilder aus Text oder Bildeingaben erzeugen, doch im Bereich der Vektorgrafiken bleibt das Feld aufgrund deren komplexer Natur zurück. Eine SVG-Datei ist ein Dokument, das Pfade, Ebenen und Hierarchie enthält, und das bedeutet, ihr Zweck ist nicht nur, etwas semantisch darzustellen, sondern sie muss auch strukturell sauber und ordentlich sein. Diese Anforderungen schaffen weit mehr Probleme, die es beim Bau eines Modells zur SVG-Generierung zu überwinden gilt. Die vorliegende Arbeit analysiert SVG-Generierungsmethoden anhand von vierzehn quelloffenen Modellen, um zu verstehen, wie interne Entwurfsentscheidungen die Struktur der Ausgabe beeinflussen. Die Analyse ist um drei Dimensionen aufgebaut: wie Modelle die Geometrie von Vektorgrafiken intern repräsentieren, welche Supervisionssignale den Generierungsprozess leiten und welche Metriken zur Bewertung der SVG nützlich sind. Repräsentation und Supervision wirken zusammen statt unabhängig voneinander und sind die zwei Hauptbestandteile des internen Designs der Modelle. Die Analyse hat gezeigt, dass dieselbe Verlustfunktion unterschiedliche strukturelle Ergebnisse erzeugen kann, je nach Repräsentationswahl, und dieselbe Repräsentation sich unter verschiedener Supervision unterschiedlich verhält. Das konsistenteste Muster über die Modelle hinweg ist ein Kompromiss zwischen visuellem Realismus und Editierbarkeit. Modelle, die visuell besonders detaillierte Ausgaben erzeugen, produzieren gleichzeitig die fragmentiertesten SVG-Dateien, während Methoden, die ihre Repräsentation einschränken oder aus Daten lernen, sauberere Ergebnisse liefern, jedoch innerhalb engerer Domänen. Aktuelle Evaluierungsmetriken (FID, CLIP, LPIPS) arbeiten ausschließlich auf gerenderten Pixeln und können strukturelle Unterschiede nicht erkennen. Um dies zu adressieren, schlägt die Arbeit ein sechsdimensionales Evaluierungsframework vor, das Pfadökonomie, geometrische Qualität, Ebenenorganisation und Editierbarkeit neben visuellen Metriken einbezieht. Die Arbeit schlägt kein neues Modell vor. Sie liefert einen Rahmen, um zu verstehen, warum aktuelle Modelle die Strukturen erzeugen, die sie erzeugen, und was sich ändern muss, damit generierte SVGs in realen Designwerkzeugen funktionieren.

\newpage
\chapter*{Acknowledgment}
I want to thank Prof. Dr. Magda Gregorova and Philipp Vaeth for supervising this thesis.
In the beginning of process, I had no clear direction and their guidance 
is what kept me focused and moving forward. 
Every session gave me something concrete to work with. 
I am also grateful for the freedom they gave me to shape the analytical approach on my own terms.

I thank Prof. Dr. Pascal Meissner, as second examiner.

Special thanks to Ismail Eyyublu, a designer at McKinsey in Baku, 
who helped me understand SVG from the perspective of someone who actually works with it. 
His explanations of what designers expect from generated vector graphics shaped 
how I thought about editability throughout this thesis.

Finally, I thank Shabnam Huseynova for her support during the thesis period.